% One-page research brief for a supervisor meeting.
% Compiles with a stock LaTeX install: pdflatex supervisor_brief.tex
\documentclass[10pt]{article}
\usepackage[a4paper,margin=0.62in]{geometry}
\usepackage{times}
\usepackage{booktabs}
\usepackage[hidelinks]{hyperref}
\usepackage{titlesec}
\usepackage{enumitem}
\titlespacing*{\section}{0pt}{0.5em}{0.2em}
\titleformat{\section}{\normalfont\normalsize\bfseries}{}{0em}{}
\setlength{\parindent}{0pt}
\setlength{\parskip}{0.32em}
\pagestyle{empty}

\begin{document}

{\LARGE\bfseries Research brief: attack-resilient embedded safety}\\[0.3em]
{\large Does core pinning isolate the safety function on a single-chip Wi-Fi SoC?}

\vspace{0.4em}
\textbf{MD Humayun} \quad M.Tech CSE (Information Security and Privacy), Adm.\ No.\ P25IS016\\
Supervisor: Dr.\ M.\ A.\ Zaveri \quad $\cdot$ \quad
\texttt{github.com/mdhumayun7/Smart-Home-Automation}

\section{The question}
Consumer intrusion alarms are built almost entirely on single-chip Wi-Fi SoCs of the ESP32 class,
where the Wi-Fi driver, the TCP/IP stack, the MQTT client and the safety application all run on the
same die, under one RTOS scheduler, sharing one interrupt controller. The standard advice is to pin
the safety task to its own core---but that advice comes from a literature studying much larger
machines, in which the attacked network interface is \emph{physically separate} from the protected
core. Does it still hold when the attacked resource is the chip's own radio driver?

\section{The gap, from the current literature}
Fourteen prior works were surveyed and individually verified against publisher records
(\texttt{docs/RELATED\_WORK.md}). The gap is narrow and real:

\begin{itemize}[leftmargin=1.2em,itemsep=0.15em,topsep=0.2em]
\item \textbf{Timing isolation under attack} is established, but only on platforms that can
separate the attacked interface from the protected core: Linux-class industrial boards
(Fischer, 2020), x86 with PREEMPT-RT and dedicated NIC queues (Meszl\'enyi, 2025), Cortex-A53 with
a hypervisor and an FPGA memory scheduler (Hoornaert, \emph{Real-Time Systems}, 2024). None of
these mechanisms---MMU, hypervisor, counter-based regulation, hardware queues---exist on a
240\,MHz MCU.
\item \textbf{Radio-layer attacks on this chip family} are cheap and demonstrated (Kremser and
Graffi, \emph{Sensors}, 2026; Cayre \emph{et al.}, WOOT'23), and deauthentication resilience has
been quantified (Carbajal, 2026)---but the dependent variable is always whether the \emph{link}
drops, never whether the device's \emph{application} still meets its deadline.
\item \textbf{ESP32 security work} measures MQTT-over-TLS cost or on-device detection accuracy.
No one treats a safety deadline as the dependent variable.
\end{itemize}

No published work measures a safety deadline on a device whose own radio driver is the attacked
resource. That is the contribution on offer.

\section{What is already established}
All of the following is committed, seeded and reproducible by a single script; CI re-runs a reduced
version of every one of them on each push.

\begin{tabular}{@{}lll@{}}
\toprule
Evidence & Scale & Result \\
\midrule
TLA+ model checking (TLC) & 2 models, 25 + 660 states & All properties hold; both mutation-tested \\
Adversarial campaign & 6{,}000 seeded trials & 0 attacker disarms, 0 invariant violations \\
Fuzzing (ASan + UBSan) & 3{,}000{,}000 iterations & 0 memory errors, 0 accepted inputs \\
\bottomrule
\end{tabular}

The second TLA+ model is the methodological point worth discussing: the network task is given
\emph{no fairness assumption}---free to stall unboundedly, flood, fail to publish or crash---and
the safety step is proved to remain enabled regardless. That is a formal statement of the isolation
claim, independent of any particular workload.

\section{What is open, and why it cannot be simulated}
Physical sensor-to-siren latency under on-chip radio stress. Simulation establishes \emph{logical}
isolation---no code path makes the safety function depend on the network---but not \emph{timing}
isolation, which is a property of the silicon, the interrupt controller and the closed-source radio
driver. The apparatus is specified and built (\texttt{experiments/hardware/}): a second ESP32 as an
independent measuring instrument, a logic probe, and collection and analysis scripts using the same
non-parametric test.

\textbf{Resources needed:} two ESP32 DevKit boards, a logic probe, and an authorised isolated or
shielded RF environment. The attack condition would be run only on our own equipment with
institutional permission; this is stated in the paper's ethics section, and no attack tooling is
published in the repository.

\section{Proposed next steps}
\begin{enumerate}[leftmargin=1.4em,itemsep=0.15em,topsep=0.2em]
\item Approve the hardware measurement campaign and the isolated-RF arrangement.
\item Run conditions A--C (idle, legitimate load, command stress) first; these need no RF isolation
and alone support a workshop submission.
\item Add condition D (radio-layer stress) once approved. This is the headline result.
\item Target venue, for discussion: an IEEE/ACM workshop on IoT or embedded security, or
\emph{Sensors}/\emph{Applied Sciences}, where comparable device-class work appears.
\end{enumerate}
\textbf{If condition D is not approved:} the paper still stands on the formal isolation result plus
conditions A--C, with D as future work. No claimed number depends on it.

\end{document}
